\chapter{Introduction} \label{cap:introduction}

The foreign exchange market (Forex or FX) trades roughly USD 7.5 trillion a day across spot, swap and derivative instruments \cite{bis2022triennial}. A market of that magnitude does not respond to price charts alone. Currency pairs react, often within seconds, to central-bank decisions, military escalation, sanctions, supply-chain shocks and political instability. Market participants able to ingest, interpret and act on such textual information faster than the average competitor retain a measurable advantage.

Trading strategies that rely exclusively on technical indicators applied to OHLC price bars \cite{murphy1999technical} have been documented to perform adequately under stable price regimes and to deteriorate when news flow, rather than price flow, drives the market. Modern Natural Language Processing (NLP) and the public release of finance-tuned language models \cite{araci2019finbert,yang2020finbert,wu2023bloomberggpt,liu2023fingpt} have substantially lowered the cost of revisiting that limitation empirically. Although the title refers to \emph{AI-driven sentiment}, the executed study tests GDELT coded-event intensity; it does not execute article-level sentiment scoring. The present work therefore investigates a hybrid trading pipeline that complements price-based signals with news-derived event features, evaluated under a strict backtesting protocol on historical data.

\section{Theme and Problem} \label{sec:theme-problem}

The theme is the integration of news-derived information into algorithmic decision-making for Forex. The driving research question is the following:

\begin{quote}
Can news-derived event features improve risk-adjusted Forex performance over a price-only baseline?
\end{quote}

\subsection*{Hypothesis} \label{sec:hypothesis}
The central premise of this study, that adding news-derived event features to a competitive price-only baseline improves risk-adjusted Forex performance, is a \textit{hypothesis}, not an established result. The evidence reviewed in Section~\ref{sec:justification} makes the premise plausible, but plausibility is not proof. The empirical apparatus developed in this work subjects that premise to a falsifiable test rather than presupposing it. The research question is therefore operationalized as a single hypothesis:
\begin{description}
  \item[\textbf{H1}] On the executed EUR/USD evaluation window, the hybrid
  strategy of Chapter~\ref{cap:methodology} produces a statistically meaningful
  improvement over the price-only baseline under the same data, execution model,
  costs and paired statistical test.
\end{description}
H1 is referenced under an explicit accept/reject criterion in
Section~\ref{sec:final-remarks} of the conclusion, derived from the evidence
reported in Chapter~\ref{cap:experimental-evaluation}. A rejection constitutes
a valid and reportable outcome: it would establish that, under this protocol
and on this pair, the event channel does not yield a statistically
meaningful improvement over the price-only baseline. Such a negative finding
is itself a contribution, and motivates stating the hypothesis in a form that
the evidence can refute.

The problem has three methodological parts. First, event data at currency-grade scale must be collected and time-aligned with FX prices \cite{leetaru2013gdelt}. Second, coded events must be transformed into numeric features that retain financially meaningful signal without leaking future information into past decisions \cite{loughran2011liability,tetlock2007giving}. Third, the combined feature space must drive entries, exits and position sizing in a way that remains auditable under regime change \cite{lopezdeprado2018advances}.

\section{Justification and Motivation} \label{sec:justification}

Three lines of reasoning motivate this investigation.

The empirical case is decades old, although most of the supporting evidence comes from equities rather than Forex. \citeonline{tetlock2007giving}, working on the US equity market, demonstrated that media tone has predictive content beyond what is contained in past prices. \citeonline{caldara2022geopolitical} constructed a geopolitical risk index from newspaper text and documented measurable real-activity and broad financial-market effects following GPR shocks. Both contributions --- Tetlock on US equities and Caldara and Iacoviello on macroeconomic and broad-market outcomes --- exemplify the dominant lines of text-based financial prediction; Forex itself has been studied predominantly through macroeconomic and microstructure variables, with comparatively little use of textual inputs. This asymmetry leaves room for an incremental contribution and motivates this work as a transposition of an equity-derived hypothesis into a Forex setting.

Two recent USP monographs sharpen this gap on both sides. On the model side, \citeonline{felix2024swing}, produced in the same ICMC/USP MBA in Artificial Intelligence and Big Data as the present work, surveys seven machine-learning classifiers for swing-trading on the IBrX-100 and explicitly identifies textual information (``new financial articles, social media and even publications from economic analysts'') as a relevant qualitative input class that the implementation nevertheless leaves out. On the asset-class side, \citeonline{vanhelden2021fx} builds a hybrid FX strategy that fuses a macro-fundamental signal (real interest-rate differential, GDP-growth differential and M2/reserves ratio) with a technical signal averaged from three moving-average crossovers, applied to twelve currency pairs. The hybrid produces an annualized Sharpe ratio above 1 on all six emerging Latin-American pairs in the sample, but \textbf{drops to 0.24 on EUR and to $-1.08$ on JPY against the USD} (\citeauthor{vanhelden2021fx}, Table~6, p.~35). The weak result on developed-market pairs motivates testing whether a different exogenous-information channel can add value to a price-only baseline.

The technological case is recent. Until 2018, financial text analytics relied predominantly on dictionary methods such as Loughran--McDonald \cite{loughran2011liability}, whose sensitivity is bounded by the contents of the lexicon. Transformer-based models \cite{vaswani2017attention,devlin2019bert} and their finance-tuned descendants \cite{araci2019finbert,yang2020finbert,wu2023bloomberggpt,liu2023fingpt} have advanced the achievable accuracy on financial sentiment well beyond the lexicon ceiling. The present work does not execute that full article-level sentiment layer, but it builds the event-based and reproducibility foundations needed for that extension.

A third motivation is methodological reusability. Conducted as an applied dissertation within an MBA in Artificial Intelligence and Big Data, the study is organized so that its pipeline can be extended. Design decisions favor permissively licensed components, explicit boundaries between perception, decision and execution, an auditable filter chain, and a perception layer decoupled from execution determinism so that model upgrades do not compromise reproducibility. The objective is a reusable methodology, not a one-off study: subsequent investigations should be able to add sentiment models, extend the asset universe or alter the evaluation protocol without reimplementing the apparatus.

\section{Objectives} \label{sec:objectives}

\subsection{General objective}

To design, implement and empirically evaluate a hybrid algorithmic-trading methodology for the Forex market, encompassing the data pipeline, the event-feature layer, the baseline and hybrid strategies and the backtesting protocol, applied to EUR/USD, and to determine, under controlled experimental conditions, whether the additional news-derived features yield a statistically meaningful improvement over a price-only baseline.

\subsection{Specific objectives}

\begin{enumerate}
    \item \textbf{Construct a reproducible data layer.} Assemble intraday EUR/USD price history and GDELT event data \cite{leetaru2013gdelt} under strict point-in-time discipline, with each completed data partition identified before it enters an experiment.

    \item \textbf{Define an event-feature layer.} Convert GDELT coded events into lagged event-intensity features aligned to the price grid. Text sentiment models, additional news corpora and the Caldara--Iacoviello GPR index \cite{caldara2022geopolitical} are retained as future work rather than treated as executed inputs.

    \item \textbf{Specify and implement baseline and hybrid strategies in a shared framework.} The baseline draws on three price-only feature families: trend, technical indicators and candlestick patterns. These are processed by per-family LightGBM sub-models whose probability outputs feed a logistic meta-learner; the meta-learner output then drives a deterministic filter chain that emits BUY, SELL or HOLD on each bar. The hybrid retains the same architecture and adds a news-derived event-intensity family as an additional sub-model in the meta-learner. Both strategies are evaluated within the same backtesting protocol and cost assumptions.

    \item \textbf{Conduct a rigorous empirical evaluation.} Apply out-of-sample risk-adjusted metrics, formal statistical tests of return-distribution differences, drift controls and ablation analyses isolating the marginal contribution of the event channel.
\end{enumerate}

\section{Document Structure} \label{sec:structure}

The remainder of this dissertation is organized as follows. Chapter~\ref{cap:theoretical-foundation} reviews the structure of the Forex market, algorithmic and quantitative trading, the NLP concepts that motivate future extensions, the literature on geopolitical risk and event data, and recent advances in finance-tuned language models. Chapter~\ref{cap:methodology} presents the executed methodology: data sources, feature engineering, model choices, backtesting protocol and evaluation criteria. Chapter~\ref{cap:experimental-evaluation} reports the experimental results, including ablations and controls that isolate the contribution of the event channel. Chapter~\ref{cap:conclusion} synthesizes the findings, discusses the limitations, and outlines directions for subsequent research.

The theoretical foundations provide context for the broader design, while the methodology is restricted to what was executed. The experimental chapter reports the baseline and hybrid strategy results, the one-year and one-trading-year protocols, the news-only arm, the adaptive-retraining frozen-policy result, the candlestick-catalog screen with its multiple-comparisons check, and the comparison with the closest prior art. Components planned but not executed are reported as future work in Chapter~\ref{cap:conclusion} rather than filled in with unsupported numbers.
